\documentclass[12pt]{article}
\usepackage[utf8]{inputenc}
\usepackage[T1]{fontenc}
\usepackage{hyperref}
\title{Waiwai Language Resources for K-12 Curriculum Development}
\author{Compiled by Claude Sonnet 5 (claude-sonnet-5) for Andr\'as Kornai}
\date{}
\begin{document}
\maketitle
\section{Introduction}
Waiwai (ISO 639-3: waw; Glottocode waiw1244) is a Cariban language of northern Brazil and southern Guyana (also reported in Suriname). Estimates are about 2,200 speakers, or 2,502 Wai-Wai people in Brazil (2014) and 170 in Guyana (2006). UNESCO lists it as vulnerable.
\section{Primary Educational Resources}
No primers or school curricula were found in search results. A SIL Brazil archive item appeared in results, but its content was not verified.
\section{Linguistic and Reference Materials}
\begin{itemize}
\item Wikipedia on the language: \url{https://en.wikipedia.org/wiki/Waiwai_language}
\item Hawkins's 1998 grammar sketch of Wai Wai is mentioned in search summaries, but no direct link was found.
\item SIL Brazil archive entry: \url{https://www.silbrazil.org/resources/archives/3517}
\item Omniglot page: \url{https://www.omniglot.com/writing/waiwai.htm}
\item Numerals dataset: \url{https://lingweb.eva.mpg.de/channumerals/Waiwai.htm}
\end{itemize}
\section{Religious and Community Resources}
Missionaries of the Unevangelized Fields Mission (the Hawkins family) developed the current orthography and translated scripture from 1949; no scripture URL was found. People page: \url{https://en.wikipedia.org/wiki/Wai-wai_people}; Indigenous Peoples in Brazil: \url{https://indios.org.br/en/Povo:Waiwai}
\section{Digital and Audio Resources}
None found.
\section{Curriculum Development Context}
There is an established orthography and a grammar, but no dictionary was found and no school materials were located.
\end{document}